\documentclass[12pt]{article}
\usepackage[T1]{fontenc}
\usepackage[utf8]{inputenc}
\usepackage{lmodern}
\usepackage{textcomp}
\usepackage{enumitem}
\usepackage{geometry}
\geometry{margin=1in}
\begin{document}
\begin{center}
Answer Key - Psychology - Graduate Level Assessment 11 \
\small (Answers are ordered exactly as in the question paper.)
\end{center}

\section*{Multiple Choice}
\begin{enumerate}[label=\arabic*.]
\item \textbf{Answer:} B
\\ \textit{Options:} A) improve external validity.; B) evaluate the impact of pretesting.; C) eliminate carryover effects.; D) increase the sample size.; E) control for confounding variables.
\\ \textit{Note:} The Solomon four-group design is specifically structured to assess the effects of pretesting on the outcomes of an experiment by including groups that either receive or do not receive a pretest, allowing for a clearer understanding of how pretesting may influence the results.
\item \textbf{Answer:} C
\\ \textit{Options:} A) Rehearsal and Forgetting; B) Sensory Input and Cognitive Load; C) Interference and Time-Decay; D) Capacity and Duration; E) Chunking and Associative Learning
\\ \textit{Note:} Interference and time-decay are the two principal explanations for the limited capacity of short-term memory, as interference arises from competing information disrupting the retention of items, while time-decay refers to the gradual loss of information over time when not actively rehearsed or encoded.
\item \textbf{Answer:} A
\\ \textit{Options:} A) A projective technique; B) A skills test; C) A reference check; D) A job sample test; E) An objective personality inventory
\\ \textit{Note:} A projective technique is considered to have the highest predictive validity in personnel selection because it reveals underlying personality traits and motivations that can predict job performance more effectively than other assessment methods.
\item \textbf{Answer:} A
\\ \textit{Options:} A) cognitive dissonance; B) positive reinforcement; C) classical conditioning; D) intrinsic motivation; E) observational learning
\\ \textit{Note:} Learned helplessness illustrates cognitive dissonance because it highlights the psychological conflict that arises when individuals perceive a lack of control over their environment, leading to a disconnect between their beliefs about personal agency and their experiences of failure.
\item \textbf{Answer:} A
\\ \textit{Options:} A) Focus solely on the influence of early childhood experiences; B) Focus mainly on the development of psychopathology over time; C) Suggest that development is a linear process without any decline; D) All of the above; E) Are mainly concerned with cognitive development
\\ \textit{Note:} The correct answer reflects a common misconception in developmental psychology, as Baltes's theory emphasizes the continuous influence of both early experiences and later life choices on development, rather than focusing exclusively on early childhood.
\end{enumerate}

\section*{True / False}
\begin{enumerate}[label=\arabic*.]
\setcounter{enumi}{5}
\item \textbf{Answer:} True
\\ \textit{Options:} A) True; B) False
\\ \textit{Note:} Indecision arises because individuals facing an avoidance-avoidance conflict struggle to choose between two undesirable options, leading to a paralysis of choice, while escalation occurs as they may prolong their decision-making process in hopes of finding a less negative outcome.
\item \textbf{Answer:} True
\\ \textit{Options:} A) True; B) False
\\ \textit{Note:} The statement is false because aggressive acts are influenced by a complex interplay of situational and psychological factors, not solely by physical characteristics.
\item \textbf{Answer:} True
\\ \textit{Options:} A) True; B) False
\\ \textit{Note:} The reticular formation and reticular activating system are primarily involved in regulating arousal and attention rather than in the processing or decussation of auditory stimuli, which mainly occurs in the brainstem and auditory pathways.
\item \textbf{Answer:} False
\\ \textit{Options:} A) True; B) False
\\ \textit{Note:} Humanistic therapists typically emphasize the importance of personal growth and self-actualization rather than attributing psychopathology primarily to severe trauma, which is more commonly associated with other therapeutic approaches such as psychodynamic or trauma-focused therapies.
\item \textbf{Answer:} True
\\ \textit{Options:} A) True; B) False
\\ \textit{Note:} The Suicide Risk Identification Tool is not widely recognized or validated in the academic literature as a standardized or established method for assessing suicide risk, which supports the assertion that it is not an established tool.
\end{enumerate}

\section*{Long Form Response}
\begin{enumerate}[label=\arabic*.]
\setcounter{enumi}{10}
\item \textbf{Answer:} Emotional content significantly enhances the effectiveness of learning verbal materials by facilitating deeper cognitive processing and retention. This phenomenon can be attributed to the interplay between the amygdala and the hippocampus, where emotional arousal activates the amygdala, thereby enhancing the encoding of memories in the hippocampus. When learners encounter verbal materials with strong emotional associations, they are more likely to engage in elaborative rehearsal, which fosters better integration of new information with existing knowledge. Additionally, emotional content can increase motivation and attention, further amplifying the learning experience. Consequently, the integration of emotional elements in educational contexts can lead to improved recall and comprehension of verbal information.
\\ \textit{Note:} correct because it identifies the role of emotional content in enhancing cognitive processing and memory retention through the interaction between the amygdala and hippocampus, along with the concept of elaborative rehearsal in learning.
\item \textbf{Answer:} Skewness refers to the asymmetry in the distribution of data, with positive skewness indicating that the tail on the right side of the distribution is longer or fatter than the left side. In analyzing the waist measurements of six randomly selected chocolate rabbits, the data reveals a positively skewed distribution. This suggests that a majority of the rabbits have smaller waist measurements, with a few outliers exhibiting significantly larger measurements. The presence of these outliers pulls the mean to the right, highlighting the skewness in the dataset. Understanding this distribution is crucial in psychological research as it can impact statistical analyses and interpretations of body image perceptions among different populations.
\\ \textit{Note:} correct because it accurately defines skewness as the asymmetry in data distributions, identifies the positive skewness in the waist measurements of the chocolate rabbits, and explains how outliers contribute to this skewness by pulling the distribution towards larger values.
\end{enumerate}

\vfill
\noindent\textit{End of Answer Key}
\end{document}